\section{Arytmetyzacja funkcji}
Numerujemy wszystkie funkcje rekurencyjne:
\begin{itemize}
\singlespacing
    \item \(\langle 0,n \rangle\) -- numer zera \(O^n\)
    \item \(\langle 1 \rangle\) -- numer następnika \(S\)
    \item \(\langle 2,n,i \rangle\) -- numer rzutowania \(I_i^n, i \leq n\)
    \item \(\langle 3,b_1,\dots,b_n,a \rangle\) -- numer złożenia:
    \[ f(\vec{x}) = g(h_1(\vec{x}), \dots, h_n(\vec{x})), \]
 gdzie \(g\) ma numer \(a\), \(h_i\) ma numer \(b_i\). 
 \item \(\langle 4,a,b \rangle\) -- numer funkcji \(f\) otrzymanej za pomocą rekursji prostej:
 \[
 \begin{cases}
     f(\vec{x},0) = g(\vec{x}) \\
     f(\vec{x},y+1) = h(\vec{x},y,f(x,y))
 \end{cases}
 \]
 \(a\) to numer funkcji \(g\), a \(b\) to numer funkcji \(h\).
 \item \(\langle 5,a \rangle\) - numer funkcji otrzymywanej przez \(\mu\)-rekursję:
 \[ f(\vec{x}) = \mu y[g(\vec{x},y)=0], \]
 gdzie \(a\) to numer funkcji \(g\).
\end{itemize}
Kodowanie nie jest suriektywne, ale możemy sprawdzić, czy kod jest poprawny, przeglądając kolejne funkcje. Dla \(f \subseteq \rec\) \: przez \(e_f\) oznaczamy najmniejszy kod \(f\).